\section{Shlukové a klasifikační úlohy strojového učení. Rozhodovací stromy, shlukovací algoritmy, asociační pravidla}


\subsection*{Klasifikační úloha}

Klasifikace je úloha učení s učitelem. Každý objekt je popsán vstupními atributy:
\[
x_i=(x_{i1},\ldots,x_{ip}),
\]
a má známou cílovou třídu:
\[
y_i\in\{1,\ldots,K\}.
\]

Cílem je naučit klasifikátor:
\[
\hat{f}(x_i)=\hat{y}_i,
\]
který dokáže přiřadit nová pozorování do správné třídy.

\paragraph{Příklady.}
\begin{itemize}
    \item zákazník odejde / neodejde,
    \item transakce je podvodná / nepodvodná,
    \item e-mail je spam / ne-spam,
    \item pacient má / nemá danou diagnózu,
    \item klient je dobrý / rizikový žadatel o úvěr.
\end{itemize}

\subsection*{Základní postup klasifikace}

\begin{enumerate}
    \item připravit data s cílovou třídou,
    \item rozdělit data na trénovací a testovací množinu,
    \item natrénovat model na trénovacích datech,
    \item ověřit model na testovacích datech,
    \item vyhodnotit metriky,
    \item případně ladit hyperparametry,
    \item použít model na nové objekty.
\end{enumerate}

\subsection*{Rozhodovací stromy}

Rozhodovací strom je klasifikační nebo regresní model složený z:
\begin{itemize}
    \item \textbf{kořene} -- první rozhodovací uzel,
    \item \textbf{vnitřních uzlů} -- testy atributů,
    \item \textbf{větví} -- výsledky testů,
    \item \textbf{listů} -- finální predikce.
\end{itemize}

Strom vytváří pravidla typu:
\[
\text{IF podmínka}_1 \wedge \text{podmínka}_2
\quad
\text{THEN třída}.
\]

\paragraph{Příklad pravidla.}
\[
\text{IF Outlook = Sunny AND Humidity = High THEN Play = No.}
\]

Výhodou stromů je, že jsou snadno interpretovatelné a lze je převést na pravidla nebo rozhodovací tabulky.

\subsection*{TDIDT}

Materiály pracují s třídou algoritmů \term{TDIDT}:
\[
\text{Top-Down Induction of Decision Trees}.
\]

Obecný postup:
\begin{enumerate}
    \item začneme všemi daty v kořeni,
    \item vybereme nejlepší atribut pro rozdělení,
    \item vytvoříme větve podle hodnot atributu nebo podle prahu,
    \item v každém poduzlu postup opakujeme,
    \item končíme, když je uzel dostatečně čistý nebo nelze dále dělit.
\end{enumerate}

Příklady algoritmů:
\begin{itemize}
    \item ID3,
    \item C4.5,
    \item C5.0,
    \item CART.
\end{itemize}

\subsection*{Čistota uzlu}

Čistota uzlu vyjadřuje, zda v uzlu převažuje jedna třída. Uzel je čistý, pokud všechna pozorování v uzlu
patří do stejné třídy.

Pokud je v uzlu směs tříd, je uzel nečistý. Algoritmus hledá takové dělení, které co nejvíce sníží nečistotu.

\subsection*{Entropie}

Entropie měří neurčitost nebo nečistotu množiny $S$.

Pro $K$ tříd:
\[
H(S)=
-\sum_{k=1}^K p_k\log_2(p_k),
\]
kde:
\[
p_k=\frac{\text{počet pozorování třídy } k}{\text{počet všech pozorování v } S}.
\]

Platí:
\begin{itemize}
    \item entropie je nulová, pokud jsou všechna pozorování ve stejné třídě,
    \item entropie je vysoká, pokud jsou třídy rovnoměrně zastoupené.
\end{itemize}

Pro dvě třídy:
\[
H(S)=
-p_+\log_2(p_+)-p_-\log_2(p_-).
\]

\subsection*{Information gain}

Information gain říká, o kolik se sníží entropie po rozdělení podle atributu $A$.

\[
Gain(S,A)
=
H(S)
-
\sum_{v\in Values(A)}
\frac{|S_v|}{|S|}
H(S_v).
\]

Kde:
\begin{itemize}
    \item $S$ je původní množina,
    \item $S_v$ je podmnožina s hodnotou atributu $A=v$,
    \item $H(S)$ je původní entropie,
    \item vážený průměr entropií poduzlů je entropie po rozdělení.
\end{itemize}

Algoritmus ID3 vybírá atribut s největším informačním ziskem:
\[
A^*
=
\underset{A}{\operatorname{argmax}}
\,
Gain(S,A).
\]

\paragraph{Intuice.}
Nejlepší atribut je takový, který po rozdělení vytvoří co nejčistší podskupiny.

\subsection*{Problém informačního zisku}

Information gain preferuje atributy s velkým počtem hodnot. Extrémní příklad je identifikátor klienta:
každý klient má unikátní ID, takže rozdělení podle ID vytvoří list pro každý záznam a trénovací data mohou být
klasifikována dokonale. Takový strom ale nebude generalizovat.

\subsection*{Gain ratio}

C4.5 proto používá normalizovanou variantu:
\[
GainRatio(S,A)
=
\frac{Gain(S,A)}
{SplitInfo(S,A)}.
\]

Kde:
\[
SplitInfo(S,A)
=
-\sum_{v\in Values(A)}
\frac{|S_v|}{|S|}
\log_2
\left(
\frac{|S_v|}{|S|}
\right).
\]

Gain ratio penalizuje atributy, které pouze mechanicky rozbijí data na mnoho malých skupin.

\subsection*{Gini index}

CART často používá Gini nečistotu:
\[
Gini(S)
=
1-\sum_{k=1}^K p_k^2.
\]

Po rozdělení podle atributu $A$:
\[
Gini(S,A)
=
\sum_{v\in Values(A)}
\frac{|S_v|}{|S|}
Gini(S_v).
\]

Vybírá se rozdělení s nejmenší nečistotou:
\[
A^*
=
\underset{A}{\operatorname{argmin}}
\,
Gini(S,A).
\]

\subsection*{Práce se spojitými atributy}

U spojitých atributů se hledá práh:
\[
x_j \leq c
\quad
\text{vs.}
\quad
x_j > c.
\]

Algoritmus typicky seřadí unikátní hodnoty atributu a zkouší možné body rozdělení. Vybere takový práh, který
maximalizuje information gain nebo gain ratio.

\paragraph{Příklad.}
\[
Temperature \leq 19
\quad
\text{vs.}
\quad
Temperature > 19.
\]

Stejný spojitý atribut se může v jedné větvi stromu objevit vícekrát, například:
\[
Temperature \leq 19
\]
a později:
\[
Temperature \leq 25.
\]

\subsection*{C4.5}

C4.5 rozšiřuje ID3. Hlavní vlastnosti:
\begin{itemize}
    \item používá gain ratio,
    \item podporuje spojité atributy,
    \item umožňuje pruning,
    \item umí pracovat s vícehodnotovými nominálními atributy,
    \item může seskupovat hodnoty atributů.
\end{itemize}

Seskupování hodnot pomáhá například u atributů s mnoha kategoriemi:
\[
\text{Monday, Tuesday, Wednesday, Thursday, Friday}
\rightarrow
\text{weekday},
\]
\[
\text{Saturday, Sunday}
\rightarrow
\text{weekend}.
\]

\subsection*{Overfitting rozhodovacích stromů}

Stromy mají silný sklon k přeučení. Pokud strom roste příliš hluboko, může se naučit šum nebo náhodné
zvláštnosti v trénovacích datech.

Typický příklad:
\begin{itemize}
    \item atribut \textit{Sunspots} náhodně odliší dva trénovací příklady,
    \item strom jej použije, protože zlepší trénovací přesnost,
    \item ale věcně nedává smysl a zhoršuje generalizaci.
\end{itemize}

\[
Accuracy_{\text{train}} \uparrow
\quad
\text{ale}
\quad
Accuracy_{\text{test}} \downarrow.
\]

\subsection*{Pruning}

Pruning omezuje složitost stromu.

\subsubsection*{Pre-pruning}

Strom se zastaví dříve, například pokud:
\begin{itemize}
    \item uzel má málo pozorování,
    \item zlepšení čistoty je malé,
    \item rozdělení není statisticky významné,
    \item byla dosažena maximální hloubka.
\end{itemize}

\subsubsection*{Post-pruning}

Nejprve se vytvoří větší strom, poté se zjednodušuje. Přespecializované podstromy se nahrazují listy.

U Reduced Error Pruning se používá validační množina. Pokud zjednodušený strom nezhorší chybu na validaci,
pruning se ponechá.

\[
Err_{\text{val}}(T')
\leq
Err_{\text{val}}(T)
\quad
\Rightarrow
\quad
\text{ponechat prořezání}.
\]

C4.5 nepoužívá Reduced Error Pruning, ale statistické prořezávání pomocí intervalů spolehlivosti, takže
nepotřebuje samostatná validační data pro pruning.

\subsection*{Převod stromu na pravidla a tabulky}

Každá cesta od kořene k listu odpovídá jednomu pravidlu:
\[
\text{IF podmínky na cestě THEN predikovaná třída}.
\]

Rozhodovací tabulka má jednotlivá pravidla jako sloupce.

Výhody:
\begin{itemize}
    \item stromy jsou dobře vizualizovatelné,
    \item pravidla jsou dobře srozumitelná,
    \item rozhodovací tabulky mohou být nejpřehlednější pro business uživatele.
\end{itemize}

\subsection*{Výhody a nevýhody rozhodovacích stromů}

Výhody:
\begin{itemize}
    \item vysoká interpretovatelnost,
    \item práce s numerickými i kategoriálními atributy,
    \item zachycení nelinearit a interakcí,
    \item snadný převod na pravidla.
\end{itemize}

Nevýhody:
\begin{itemize}
    \item sklon k přeučení,
    \item nestabilita -- malá změna dat může změnit strom,
    \item preferování některých typů atributů,
    \item horší predikční výkon než ensemble metody.
\end{itemize}

\subsection*{Shlukování}

Shlukování je úloha učení bez učitele. Nemáme cílovou proměnnou, ale chceme najít skupiny podobných objektů.

Cíl:
\[
\text{minimalizovat vnitroshlukovou vzdálenost}
\]
a zároveň:
\[
\text{maximalizovat mezishlukovou vzdálenost}.
\]

Shluk by měl být:
\begin{itemize}
    \item interně homogenní,
    \item externě odlišný,
    \item pokud možno izolovaný.
\end{itemize}

\subsection*{Míry vzdálenosti a podobnosti}

Pro numerická data se často používá eukleidovská vzdálenost:
\[
d(x,z)=
\sqrt{
\sum_{j=1}^p
(x_j-z_j)^2
}.
\]

Manhattanská vzdálenost:
\[
d(x,z)=
\sum_{j=1}^p
|x_j-z_j|.
\]

Hammingova vzdálenost pro kategoriální nebo binární atributy:
\[
d(x,z)=
\sum_{j=1}^p
I(x_j\neq z_j).
\]

Jaccardova podobnost pro množiny:
\[
J(A,B)=
\frac{|A\cap B|}
{|A\cup B|}.
\]

\paragraph{Důležité.}
U shlukování je často nutná standardizace proměnných. Jinak proměnná s velkým měřítkem dominuje vzdálenosti.

\[
z_{ij}=
\frac{x_{ij}-\bar{x}_j}{s_j}.
\]

\subsection*{Centroid}

Centroid je střed shluku:
\[
\mu_k
=
\frac{1}{|C_k|}
\sum_{x_i\in C_k}
x_i.
\]

Nemusí odpovídat skutečnému pozorování. Slouží jako reprezentant shluku.

\subsection*{K-means}

K-means je nehierarchická shlukovací metoda. Vstupem je počet shluků $k$.

Cílová funkce:
\[
\min_{C_1,\ldots,C_k}
\sum_{r=1}^k
\sum_{x_i\in C_r}
\|x_i-\mu_r\|^2.
\]

Postup:
\begin{enumerate}
    \item zvolíme počet shluků $k$,
    \item inicializujeme centroidy,
    \item přiřadíme každý objekt k nejbližšímu centroidu,
    \item přepočítáme centroidy,
    \item opakujeme, dokud se přiřazení nemění nebo není dosažen limit iterací.
\end{enumerate}

\subsubsection*{Parametry k-means}

\begin{itemize}
    \item počet shluků $k$,
    \item metrika vzdálenosti,
    \item inicializace centroidů,
    \item počet náhodných startů,
    \item maximální počet iterací,
    \item zastavovací kritérium.
\end{itemize}

\subsubsection*{Výhody a nevýhody k-means}

Výhody:
\begin{itemize}
    \item jednoduchý a rychlý algoritmus,
    \item dobře škáluje na větší data,
    \item snadná interpretace přes centroidy.
\end{itemize}

Nevýhody:
\begin{itemize}
    \item nutnost předem zadat $k$,
    \item citlivost na inicializaci,
    \item citlivost na outliery,
    \item předpokládá spíše kulovité shluky,
    \item typicky vyžaduje numerická data.
\end{itemize}

\subsection*{Volba počtu shluků}

\subsubsection*{Elbow method}

Sledujeme vnitroshlukový součet čtverců:
\[
WCSS(k)
=
\sum_{r=1}^k
\sum_{x_i\in C_r}
\|x_i-\mu_r\|^2.
\]

S rostoucím $k$ hodnota klesá. Hledáme bod zlomu, po kterém přidání dalšího shluku už příliš nezlepší
výsledek.

\subsubsection*{Silhouette score}

Pro pozorování $i$:
\[
s_i=
\frac{b_i-a_i}{\max(a_i,b_i)}.
\]

Kde:
\begin{itemize}
    \item $a_i$ je průměrná vzdálenost objektu od bodů ve stejném shluku,
    \item $b_i$ je nejmenší průměrná vzdálenost k bodům v jiném shluku.
\end{itemize}

Platí:
\[
-1\leq s_i \leq 1.
\]

Vyšší hodnota znamená lepší přiřazení objektu ke shluku.

\subsection*{Hierarchické shlukování}

Hierarchické shlukování vytváří hierarchii shluků. Výsledkem je dendrogram.

\subsubsection*{Aglomerativní shlukování}

Postup zdola nahoru:
\begin{enumerate}
    \item každý objekt je samostatný shluk,
    \item najdeme dva nejbližší shluky,
    \item sloučíme je,
    \item aktualizujeme vzdálenosti,
    \item opakujeme až do jednoho shluku nebo požadovaného počtu shluků.
\end{enumerate}

\subsubsection*{Divizní shlukování}

Postup shora dolů:
\begin{enumerate}
    \item začneme jedním shlukem obsahujícím všechna pozorování,
    \item postupně jej dělíme,
    \item končíme při požadovaném počtu shluků.
\end{enumerate}

\subsection*{Vzdálenosti mezi shluky}

Single linkage:
\[
D(C_g,C_h)
=
\min_{x_i\in C_g,\,x_j\in C_h}
d(x_i,x_j).
\]

Complete linkage:
\[
D(C_g,C_h)
=
\max_{x_i\in C_g,\,x_j\in C_h}
d(x_i,x_j).
\]

Average linkage:
\[
D(C_g,C_h)
=
\frac{1}{|C_g||C_h|}
\sum_{x_i\in C_g}
\sum_{x_j\in C_h}
d(x_i,x_j).
\]

Centroid linkage:
\[
D(C_g,C_h)=d(\mu_g,\mu_h).
\]

Wardova metoda minimalizuje nárůst vnitroshlukového součtu čtverců při sloučení shluků.

\subsection*{Asociační pravidla}

Asociační pravidla hledají vztahy mezi položkami nebo hodnotami atributů. Typické použití je analýza
nákupního košíku.

Pravidlo:
\[
Ant \Rightarrow Suc.
\]

Například:
\[
\{bread,butter\}\Rightarrow\{milk\}.
\]

Antecedent je levá strana pravidla, sukcedent je pravá strana.

\subsection*{Transakční data}

Transakce je množina položek:
\[
T_i=\{milk,bread,butter\}.
\]

Dataset je množina transakcí:
\[
D=\{T_1,\ldots,T_n\}.
\]

Itemset je množina položek:
\[
X=\{bread,butter\}.
\]

\subsection*{Kontingenční tabulka pravidla}

Pro pravidlo:
\[
Ant \Rightarrow Suc
\]
definujeme:

\[
\begin{array}{c|cc}
 & Suc & \neg Suc\\
\hline
Ant & a & b\\
\neg Ant & c & d
\end{array}
\]

Kde:
\begin{itemize}
    \item $a$ splňuje antecedent i sukcedent,
    \item $b$ splňuje antecedent, ale ne sukcedent,
    \item $c$ nesplňuje antecedent, ale splňuje sukcedent,
    \item $d$ nesplňuje ani antecedent, ani sukcedent.
\end{itemize}

\subsection*{Základní charakteristiky pravidel}

Podpora:
\[
support(Ant\Rightarrow Suc)
=
\frac{a}{a+b+c+d}.
\]

Někdy se podpora počítá absolutně jako počet případů $a$.

Confidence:
\[
confidence(Ant\Rightarrow Suc)
=
P(Suc|Ant)
=
\frac{a}{a+b}.
\]

Coverage:
\[
coverage(Ant\Rightarrow Suc)
=
P(Ant|Suc)
=
\frac{a}{a+c}.
\]

Lift:
\[
lift(Ant\Rightarrow Suc)
=
\frac{P(Suc|Ant)}{P(Suc)}
=
\frac{a/(a+b)}{(a+c)/(a+b+c+d)}.
\]

Interpretace liftu:
\begin{itemize}
    \item $lift>1$ znamená pozitivní asociaci,
    \item $lift=1$ znamená nezávislost,
    \item $lift<1$ znamená negativní asociaci.
\end{itemize}

\subsection*{Apriori algoritmus}

Apriori hledá časté množiny položek a z nich generuje pravidla.

Klíčová vlastnost:
\[
\text{Pokud je itemset častý, všechny jeho podmnožiny jsou také časté.}
\]

Ekvivalentně:
\[
\text{Pokud itemset není častý, žádná jeho nadmnožina nemůže být častá.}
\]

To umožňuje výrazně omezit počet kandidátů.

\subsubsection*{Postup Apriori}

\begin{enumerate}
    \item najdi časté jedno-položkové množiny $F_1$,
    \item z $F_k$ vytvoř kandidáty $C_{k+1}$,
    \item odstraň kandidáty, jejichž některá podmnožina není častá,
    \item spočítej podporu kandidátů v databázi,
    \item ponech kandidáty s podporou alespoň $min\_support$,
    \item pokračuj, dokud nevzniknou další časté množiny.
\end{enumerate}

\[
F_{k+1}=
\{X\in C_{k+1}: support(X)\geq min\_support\}.
\]

Z častého itemsetu $X$ lze generovat pravidla:
\[
A\Rightarrow X\setminus A,
\quad
A\subset X.
\]

Pravidlo se ponechá, pokud:
\[
confidence(A\Rightarrow X\setminus A)\geq min\_confidence.
\]

\subsection*{GUHA a LISp-Miner}

GUHA je česká tradice dobývání znalostí z databází. Kombinuje generování hypotéz s jejich ověřováním.
LISp-Miner je systém navazující na GUHA. Procedura 4FT-Miner hledá asociační vztahy v datech pomocí
kontingenčních tabulek a zadaných kritérií zajímavosti.

\subsection*{CBA}

CBA, Classification Based on Associations, používá asociační pravidla pro klasifikaci. Hledají se pravidla,
kde pravá strana obsahuje třídu:
\[
Ant \Rightarrow class.
\]

Z pravidel se vytvoří klasifikátor. Pravidla se typicky řadí podle spolehlivosti, podpory a jednoduchosti.

\subsection*{Srovnání klasifikace, shlukování a asociačních pravidel}

\begin{itemize}
    \item \textbf{Klasifikace:}
    máme cílovou proměnnou a učíme predikovat třídu.
    
    \item \textbf{Shlukování:}
    nemáme cílovou proměnnou a hledáme přirozené skupiny objektů.
    
    \item \textbf{Asociační pravidla:}
    nehledáme třídy ani shluky, ale vztahy mezi položkami nebo atributy.
\end{itemize}

\subsection*{Shrnutí otázky}

Klasifikační úlohy přiřazují objekty do předem známých tříd. Rozhodovací stromy jsou interpretovatelné
klasifikátory, které postupně dělí data podle atributů. ID3 používá entropii a information gain, C4.5 přidává
gain ratio, podporu spojitých atributů a pruning. Stromy jsou srozumitelné, ale snadno se přeučí.

Shlukování je učení bez učitele. Cílem je najít homogenní a dobře oddělené skupiny. K-means minimalizuje
vzdálenost objektů od centroidů, hierarchické metody vytvářejí dendrogram.

Asociační pravidla hledají souvislosti typu:
\[
Ant \Rightarrow Suc.
\]
Hodnotí se pomocí podpory, confidence a liftu. Apriori používá vlastnost, že nadmnožina nečasté množiny
nemůže být častá.

\newpage